\documentclass[oneside]{scrbook}

% \usepackage[utf8]{inputenc}
\usepackage[T1]{fontenc}
\usepackage{version}
% \usepackage[latin1]{inputenc}
\usepackage[italian]{babel}
\usepackage{algpseudocode}
\usepackage{amsmath}
\usepackage{booktabs}
\usepackage{color}
\usepackage{amsfonts}
\usepackage{url}
\usepackage{algorithm}
\usepackage{chngcntr}
\usepackage[answerdelayed]{exercise}
\usepackage{graphicx}
\usepackage{caption}
\usepackage{listings}
\usepackage{tikz}
\usepackage{pgf}
\usepackage{colortbl}
\usepackage{listingsutf8}
\usepackage[pdftitle = {Dispense di esercizi}, pdfauthor = {Gianmarco Tedeschi e Alberto Marchesi}, pdfsubject = {Fondamenti di Informatica ($082746$ - A.A.~$2025-2026$)}, pdfcreator = {Gianmarco Tedeschi e Alberto Marchesi}, pdfproducer = {-}]{hyperref}
% \usepackage[automark,nouppercase]{scrpage2}
\usepackage{mdframed}
\usepackage{scrlayer-scrpage} % Necessary package for page styles

%\pagestyle{scrheadings}

%\usepackage{titlesec}
%\newcommand{\sectionbreak}{\clearpage}
\counterwithin{Exercise}{chapter}
\counterwithin{Answer}{chapter}

\algrenewcommand\algorithmicthen{\textbf{Allora}}
\algrenewcommand\algorithmicif{\textbf{Se}}
\algrenewcommand\algorithmicelse{\textbf{Altrimenti}}
\algrenewcommand\algorithmicend{\textbf{Chiudi}}
\algrenewcommand\algorithmicwhile{\textbf{Finch\'e}}
\algrenewcommand\algorithmicdo{\textbf{Esegui}}
\algrenewcommand{\algorithmiccomment}[1]{\hskip3em\texttt{//#1}}

\setlength{\parindent}{0in}
\setlength{\parskip}{10pt}

\newmdenv[%
        outerlinewidth=2,%
        roundcorner=10pt,%
        leftmargin=0,%
        rightmargin=0,%
        outerlinecolor=black,
        innertopmargin=\topskip,%
        splittopskip=\topskip,%
]{algo}


\ohead{\headmark}
\ihead{}
\chead{}
\ifoot{A.A.~$2025-2026$}
\cfoot{Fondamenti di Informatica (082746) - G. Tedeschi}
\ofoot{Pag.~\pagemark}
\setheadsepline{.4pt}
\setfootsepline{.4pt}

\definecolor{lbcolor}{rgb}{0.9,0.9,0.9}

\lstdefinestyle{cstyle}{%
  inputencoding=utf8/latin1,
  breaklines=true,
  breakatwhitespace=false,
  basicstyle=\small\ttfamily,
  language=c,
  frame=single,
  showstringspaces=false
}

\lstdefinestyle{cstylesmall}{%
  inputencoding=utf8/latin1,
  breaklines=true,
  breakatwhitespace=false,
  basicstyle=\scriptsize\ttfamily,
  language=c,
  frame=single,
  showstringspaces=false
}

\lstdefinestyle{console}{%
  inputencoding=utf8/latin1,
  breaklines=true,
  breakatwhitespace=false,
  basicstyle=\scriptsize\ttfamily,
  showstringspaces=false,
  backgroundcolor=\color{lbcolor}
}


\lstnewenvironment{clang}{\lstset{style=cstyle}}{}
\lstnewenvironment{clangsmall}{\lstset{style=cstylesmall}}{}

\addtocounter{chapter}{7}
\begin{document}

\chapter{Memoria Dinamica e Liste}

Questa dispensa approfondisce la gestione dinamica della memoria in C e l'implementazione delle liste concatenate, strutture dati fondamentali quando la quantità di dati da gestire non è nota a priori.

\section{Allocazione Dinamica della Memoria}

In un programma C, la memoria è partizionata in diverse zone. Le variabili statiche e automatiche (quelle locali alle funzioni) risiedono nello \textbf{Stack}. Tuttavia, spesso non si conosce il numero esatto di dati da allocare fino al momento dell'esecuzione (run-time).


Per gestire questi casi, si utilizza lo \textbf{Heap} (Mucchio), un'area di memoria dedicata alle variabili dinamiche. A differenza delle variabili automatiche, quelle dinamiche:
\begin{itemize}
    \item Vengono allocate e de-allocate esplicitamente dal programmatore.
    \item Sono accessibili esclusivamente tramite puntatori.
\end{itemize}

\paragraph{Funzioni \texttt{malloc} e \texttt{free}}

La libreria standard \texttt{stdlib.h} fornisce le funzioni necessarie per l'allocazione dinamica della memoria.

\begin{itemize}
    \item \textbf{\texttt{malloc} (Memory Allocation)}: Alloca un blocco di byte contigui nello Heap. Restituisce un puntatore generico \texttt{void*} all'indirizzo base del blocco allocato. Se lo spazio non è sufficiente, restituisce \texttt{NULL}.
    \item \textbf{\texttt{free}}: De-alloca un blocco di memoria precedentemente riservato, rendendolo disponibile per nuovi utilizzi.
\end{itemize}

Esempio di allocazione di un intero:
\begin{lstlisting}[style=cstyle]
int *ptr;
// Allocazione nello Heap: sizeof calcola i byte necessari
ptr = (int*) malloc(sizeof(int)); 

if (ptr != NULL) {
    *ptr = 10;      // Utilizzo della variabile dinamica
    printf("%d", *ptr);
    free(ptr);      // Rilascio della memoria
}
\end{lstlisting}

\section{Liste Concatenate}

Le liste sono strutture dati dinamiche che possono crescere e decrescere durante l'esecuzione. A differenza degli array, gli elementi di una lista non sono necessariamente contigui in memoria; sono sparpagliati nello Heap e collegati logicamente tramite puntatori.


Ogni elemento (o nodo) contiene due parti fondamentali:
\begin{enumerate}
    \item \textbf{Dato}: l'informazione utile.
    \item \textbf{Link}: un puntatore al nodo successivo.
\end{enumerate}
L'ultimo nodo della lista ha il puntatore impostato a \texttt{NULL} per indicare la fine della sequenza.

\subsection{Definizione del Tipo}
La definizione della struttura è ricorsiva (o auto-referenziale), poiché contiene un puntatore a una struttura dello stesso tipo.

\begin{lstlisting}[style=cstyle]
typedef struct EL {
    int dato;           // Informazione (es. intero)
    struct EL * prox;   // Link al prossimo elemento
} ElemLista;

typedef ElemLista * ListaDiElem; // Puntatore alla testa
\end{lstlisting}

\section{Operazioni sulle Liste}

Le operazioni vengono implementate tramite funzioni che manipolano i puntatori. Spesso queste funzioni restituiscono il nuovo puntatore alla testa della lista, poiché le operazioni di inserimento o cancellazione potrebbero modificarlo.

\subsection{Inizializzazione e Calcolo Dimensione}

Una lista vuota è rappresentata semplicemente dal valore \texttt{NULL}.

\begin{lstlisting}[style=cstyle]
ListaDiElem inizializza() {
    return NULL; // Lista vuota
}

// Calcolo della dimensione (Versione Iterativa)
int dimensione(ListaDiElem lista) {
    int count = 0;
    while (lista != NULL) {   // Scorre fino alla fine
        lista = lista->prox;  // Avanza al prossimo nodo
        count++;
    }
    return count; // Restituisce il conteggio
}
\end{lstlisting}

Esiste anche una versione ricorsiva, basata sull'induzione: la lunghezza è 0 se la lista è vuota, altrimenti è 1 + la lunghezza del resto della lista.

\begin{lstlisting}[style=cstyle]
int dimensione(ListaDiElem lista) {
    // 1. Caso base: siamo alla fine della lista (o vuota)
    if (lista == NULL) {
        return 0;
    } else {
        // 2. Passo ricorsivo: 1 (me stesso) + la dimensione del resto
        return 1 + dimensione(lista->prox);
    }
}
\end{lstlisting}

\subsection{Inserimento}

L'inserimento richiede l'allocazione dinamica di un nuovo nodo e l'aggiornamento dei collegamenti.

\paragraph{Inserimento in Testa}
È l'operazione più efficiente ($O(1)$). Il nuovo nodo punta alla vecchia testa, e diventa la nuova testa.


\begin{lstlisting}[style=cstyle]
ListaDiElem inserisciTesta(ListaDiElem lista, int elem) {
    ListaDiElem punt;
    punt = (ListaDiElem) malloc(sizeof(ElemLista)); // Alloca
    if(punt == NULL) {
        // Gestione errore di allocazione
        return lista; // Ritorna la lista originale
    }

    punt->dato = elem;
    punt->prox = lista; // Collega al vecchio primo elemento
    
    return punt;        // Ritorna il nuovo primo elemento
}
\end{lstlisting}

\paragraph{Inserimento in Coda}
Richiede la scansione di tutta la lista per trovare l'ultimo elemento e agganciarvi il nuovo nodo.

\begin{lstlisting}[style=cstyle]
ListaDiElem inserisciCoda(ListaDiElem lista, int elem) {
    ListaDiElem punt, cur = lista;
    punt = (ListaDiElem) malloc(sizeof(ElemLista));
    punt->dato = elem;
    punt->prox = NULL; // Sarà l'ultimo

    if (lista == NULL) {
        return punt; // Se vuota, il nuovo nodo è la testa
    } else {
        while (cur->prox != NULL) // Cerca l'ultimo
            cur = cur->prox;
        cur->prox = punt; // Aggancia
        return lista;
    }
}
\end{lstlisting}

\paragraph{Inserimento Ordinato}
Mantiene la lista ordinata cercando la posizione corretta. Richiede due puntatori: \texttt{puntCor} (corrente) e \texttt{puntPrec} (precedente) per gestire l'inserimento nel mezzo.

\begin{lstlisting}[style=cstyle]
ListaDiElem inserisciOrd(ListaDiElem lista, int elem) {
    ListaDiElem punt, puntCor = lista, puntPrec = NULL;
    
    // Trova la posizione
    while (puntCor != NULL && elem > puntCor->dato) {
        puntPrec = puntCor;
        puntCor = puntCor->prox;
    }
    punt = (ListaDiElem) malloc(sizeof(ElemLista));
    punt->dato = elem;
    if (puntPrec != NULL) { // Inserimento interno o in coda
        punt->prox = puntCor;
        puntPrec->prox = punt;
        return lista;
    } else { // Inserimento in testa
        punt->prox = lista;
        return punt;
    }
}
\end{lstlisting}

\subsection{Cancellazione}

Per rimuovere un nodo, bisogna "scavalcarlo" aggiornando il puntatore del nodo precedente, e poi liberare la memoria con \texttt{free}.


Ecco un esempio di funzione ricorsiva per cancellare la prima occorrenza di un elemento:

\begin{lstlisting}[style=cstyle]
ListaDiElem cancella(ListaDiElem lista, int elem) {
    ListaDiElem puntTemp;
    if (lista != NULL) {
        if (lista->dato == elem) {
            // Elemento trovato: salva il prossimo e libera il corrente
            puntTemp = lista->prox;
            free(lista);
            return puntTemp; // Il prossimo diventa la testa (o il link corrente)
        } else {
            // Ricorsione sul resto della lista
            lista->prox = cancella(lista->prox, elem);
        }
    }
    return lista;
}
\end{lstlisting}

\subsection{De-allocazione Completa}
Quando una lista non serve più, è fondamentale liberare tutta la memoria allocata scorrendo la lista e chiamando \texttt{free} su ogni nodo.

\begin{lstlisting}[style=cstyle]
void eliminaLista(ListaDiElem lista) {
    ListaDiElem temp;
    while (lista != NULL) {
        temp = lista->prox; // Salva il riferimento al prossimo
        free(lista);        // Libera il corrente
        lista = temp;       // Avanza
    }
}
\end{lstlisting}

\section{Esercizi}

% Facili
\begin{Exercise}[label={ex:minimo_pari}]
    Data una lista di interi, scrivere una funzione che cerca il minimo elemento pari e 
    restituisce il puntatore all'area di memoria contenente il nodo associato a quel valore.
\end{Exercise}

\begin{Answer}[ref={ex:minimo_pari}]
  	\lstinputlisting[style=cstylesmall]{src/minimo_pari.c}
\end{Answer}
\begin{Exercise}[label={ex:somma_multipli}]
    Si consideri una lista di interi. Scrivere una funzione che riceve come parametro la lista ed un intero M
    e che restituisce la somma dei soli valori della lista che sono multipli di M. Se la lista è vuota, la funzione
    deve restituire -1.
\end{Exercise}

\begin{Answer}[ref={ex:somma_multipli}]
  	\lstinputlisting[style=cstylesmall]{src/somma_multipli.c}
\end{Answer}
\begin{Exercise}[label={ex:duplicati}]
    Scrivere una funzione che elimina gli elementi duplicati da una lista di interi.
    Si implementi sia la versione iterativa che quella ricorsiva.
\end{Exercise}

\begin{Answer}[ref={ex:duplicati}]
  	\lstinputlisting[style=cstylesmall]{src/duplicati.c}
\end{Answer}
\begin{Exercise}[label={ex:monotona}]
    Scrivere una funzione che, ricevuta in ingresso una lista di interi, verifichi
    se l'andamento della lista è monotono crescente (ogni elemento è strettamente maggiore al precedente).
\end{Exercise}

\begin{Answer}[ref={ex:monotona}]
  	\lstinputlisting[style=cstylesmall]{src/monotona.c}
\end{Answer}
\begin{Exercise}[label={ex:inverti}]
    Si consideri una lista di interi concatenata.
    Scrivere una funzione che inverte la lista. Due soluzioni sono possibili:
    \begin{itemize}
        \item la funzione alloca in memoria una nuova lista (invertita) e poi (se richiesto) elimina
        la lista originale;
        \item la funzione riusa i nodi della lista per crearne una nuova invertita.
    \end{itemize}
\end{Exercise}

\begin{Answer}[ref={ex:inverti}]
  	\lstinputlisting[style=cstylesmall]{src/inverti.c}
\end{Answer}
\begin{Exercise}[label={ex:bigger}]
    Definire una funzione che, data una lista opportunamente definita ed un intero n, 
    inserisce n dopo l’ultimo elemento della lista maggiore di n. Se la lista non contiene
    alcun elemento maggiore di n, la procedura lascia la lista inalterata.
\end{Exercise}

\begin{Answer}[ref={ex:bigger}]
  	\lstinputlisting[style=cstylesmall]{src/bigger.c}
\end{Answer}

% Difficili
\begin{Exercise}[label={ex:picchi}]
    Si consideri una lista di interi. 
    Definiamo picchi della lista quei valori che sono preceduti e seguiti da valori
    più piccoli della loro metà. Il primo e l'ultimo elemento della lista non possono essere picchi.
    Scrivere una funzione che data una lista restituisce il numero di picchi della lista.
\end{Exercise}

\begin{Answer}[ref={ex:picchi}]
  	\lstinputlisting[style=cstylesmall]{src/picchi.c}
\end{Answer}
\begin{Exercise}[label={ex:compressa}]
    Si considerino le seguenti definizioni:
    \begin{lstlisting}[style=cstylesmall]
        // Lista di interi
        typedef struct el{
            int dato;
            struct el* prox;
        } Node;
        typedef Node* Lista;

        // Lista compressa di interi
        typedef struct c_el{
            int dato, count;
            struct c_el* prox;
        } CNode;
        typedef CNode* ListaCompressa;
    \end{lstlisting}
    Le liste compresse rappresentano le classiche liste con nodi in cui appare il valore e il numero di volte che appare nella lista originale.
    
    \textbf{Esempio:} la lista \texttt{[3,3,3,5,5,2,2,2,2]} viene rappresentata come \texttt{[(3,3),(5,2),(2,4)]}.  
    Scrivere una funzione che prende una lista e costruisce la sua rappresentazione compressa.
\end{Exercise}

\begin{Answer}[ref={ex:compressa}]
  	\lstinputlisting[style=cstylesmall]{src/compressa.c}
\end{Answer}
\begin{Exercise}[label={ex:simmetrica}]
    Si scriva una funzione che calcola la differenza simmetrica degli elementi di due liste ordinate 
    in senso crescente e prive di duplicati, restituendola come una nuova lista (allocata allo scopo), 
    anch’essa ordinata. La differenza simmetrica è costituita dagli elementi che appartengono a una delle due liste ma non all’altra
    lista (contiene cioè tutti gli elementi che non sono in comune alle due liste).
\end{Exercise}

\begin{Answer}[ref={ex:simmetrica}]
  	\lstinputlisting[style=cstylesmall]{src/simmetrica.c}
\end{Answer}
\begin{Exercise}[label={ex:trova_k}]
    Scrivere una funzione che, ricevendo in ingresso una lista dinamica e un intero 
    $k > 0$, individui il nodo che si trova k posizioni dopo il nodo centrale 
    e ne restituisca il puntatore al chiamante. 
    Se tale nodo non esiste, ad esempio a causa di un’insufficiente numero di nodi 
    nella lista, la funzione deve restituire il valore NULL.
    
    \textbf{BONUS}: scrivere una soluzione che prevede una sola scansione 
    della lista (una scansione intera più una scansione anche solo parziale 
    è considerata equivalente a due scansioni).
\end{Exercise}

\begin{Answer}[ref={ex:trova_k}]
  	\lstinputlisting[style=cstylesmall]{src/trova_k.c}
\end{Answer}
\begin{Exercise}[label={ex:ondulatorio}]
    Si scriva una funzione che verifichi se l'andamento di una lista è \emph{ondulatorio}, cioè se non 
    capita mai che tre numeri consecutivi siano in ordine crescente o decrescente. La funzione restituisce 1
    se l'andamento è ondulatorio, 0 altrimenti.
\end{Exercise}

\begin{Answer}[ref={ex:ondulatorio}]
  	\lstinputlisting[style=cstylesmall]{src/ondulatorio.c}
\end{Answer}
\clearpage
\begin{Exercise}[label={ex:flatten}]
    Si definisca una struttura dati per una lista concatenata in cui ogni nodo è 
    costituito da un valore intero, un puntatore al nodo successivo e un puntatore opzionale (\texttt{child}) 
    che identifica la testa di una sotto-lista gerarchica. 
    Si scriva un programma che costruisce tale struttura e implementa una funzione di inserimento e di stampa che gestisca
    la gerarchizzazione della struttura ed una funzione di "appiattimento".
    (flatten) che linearizza la gerarchia: ogni sotto-lista puntata da \texttt{child} deve essere inserita nella lista principale 
    immediatamente dopo il nodo genitore, modificando i puntatori esistenti senza allocare nuovi nodi. 
    Il programma visualizzi a video la sequenza lineare risultante.
    
    \textbf{Esempio:} data la lista gerarchica
    \begin{verbatim}
        1 -> 2 -> 5
             |
             3 -> 4
                  |
                  10
    \end{verbatim}
    la lista appiattita sarà:
    \begin{verbatim}
        1 -> 2 -> 3 -> 4 -> 10 -> 5
    \end{verbatim}
\end{Exercise}

\begin{Answer}[ref={ex:flatten}]
  	\lstinputlisting[style=cstylesmall]{src/flatten.c}
\end{Answer}

\clearpage
\section*{Soluzioni}
\shipoutAnswer

\clearpage
\section*{Libreria Liste}
\lstinputlisting[style=cstylesmall]{src/lista.c}

\end{document}